\section{September 23, 2026}
% Based on pages 1--4 of 551-0923.pdf.

We explain why quadrature precision determines the endpoint order of
a collocation method, then construct Gauss Runge--Kutta methods.
Exactness of Gaussian quadrature also gives preservation of quadratic
first integrals and a direct proof of \(A\)-stability.
We retain \(\tau=\tau_n\), the stage slopes \(k_i\), and the RK matrix
\(\mathsf A\); \(A\) is reserved for an ODE matrix.

\subsection{Collocation and quadrature error}

Recall that the collocation polynomial \(u\in(\mathcal P_s)^d\)
satisfies
\begin{align*}
 u(t_n)&=x_n,\\
 u'(t_n+c_i\tau)&=f\bigl(t_n+c_i\tau,u(t_n+c_i\tau)\bigr),
       \qquad i=1,\ldots,s,\\
 x_{n+1}&=u(t_n+\tau).
\end{align*}
The associated RK coefficients are
\[
 a_{ij}=\int_0^{c_i}\ell_j(\theta)\,d\theta,
 \qquad b_j=\int_0^1\ell_j(\theta)\,d\theta,
\]
where \(\ell_j\) is the Lagrange basis from
\eqref{eq:lagrange-basis-sep-21}.

On the physical time interval, define the quadrature error functional
\begin{equation}\label{eq:physical-quadrature-error-sep-23}
 E_{n,\tau}[\varphi]
  =\int_{t_n}^{t_n+\tau}\varphi(t)\,dt
       -\tau\sum_{j=1}^s b_j\varphi(t_n+c_j\tau).
\end{equation}
This is a linear functional. As on September 21, \(E\) denotes
quadrature error, while \(R(z)\) denotes a stability function.

\begin{lemma}[Scaling of quadrature error]
 Suppose the quadrature rule has degree of precision \(p\). For
 \(\varphi\in C^{p+1}([t_n,t_n+\tau])\),
 \[
  |E_{n,\tau}[\varphi]|
   \leq C\tau^{p+2}
      \max_{t_n\leq t\leq t_n+\tau}|\varphi^{(p+1)}(t)|,
 \]
 where \(C\) depends only on the fixed nodes and weights.
\end{lemma}

\begin{proof}
 Expand \(\varphi(t_n+\tau\theta)\) through degree \(p\):
 \[
  \varphi(t_n+\tau\theta)
   =\sum_{m=0}^p\frac{\varphi^{(m)}(t_n)}{m!}
                     \tau^m\theta^m+r(\theta),
 \]
 with
 \[
  |r(\theta)|\leq
   \frac{\tau^{p+1}}{(p+1)!}
   \max|\varphi^{(p+1)}|,
   \qquad 0\leq\theta\leq1.
 \]
 Exactness cancels every displayed polynomial term. The remaining
 integral and weighted sum are bounded by
 \(\tau(1+\sum_j|b_j|)\max|r|\), proving the estimate.
\end{proof}

The same estimate holds componentwise for vector-valued functions.
If the integrand itself depends on \(\tau\), the derivative bound
must be uniform as \(\tau\to0\).

\subsection{The collocation defect and variation of constants}

For the local error analysis, let \(x(t)\) be the exact solution
started from \(x(t_n)=x_n\). Define the defect
\begin{equation}\label{eq:collocation-defect-sep-23}
 h(t)=u'(t)-f(t,u(t)).
\end{equation}
Then the polynomial solves a perturbed differential equation,
\[
 u'=f(t,u)+h(t),\qquad u(t_n)=x_n,
\]
and the collocation conditions say exactly that
\begin{equation}\label{eq:defect-zeros-sep-23}
 h(t_n+c_i\tau)=0,\qquad i=1,\ldots,s.
\end{equation}

For the linear problem \(x'=Ax\), variation of constants gives
\[
 u(t)=e^{(t-t_n)A}x_n+
          \int_{t_n}^t e^{(t-\eta)A}h(\eta)\,d\eta.
\]
Thus, writing \(T=t_n+\tau\),
\begin{equation}\label{eq:linear-defect-error-sep-23}
 x(T)-u(T)=-\int_{t_n}^T e^{(T-t)A}h(t)\,dt.
\end{equation}
The minus sign follows from defining \(h=u'-f(t,u)\) and taking
the local error to be exact minus numerical, as in
\eqref{eq:local-error-aug-31}.

For a nonlinear equation, the analogous identity uses the exact flow
\(\Phi^{T,t}(y)\), namely the solution at time \(T\) started from
\(y\) at time \(t\). Put
\[
 M(T,t)=D_y\Phi^{T,t}(u(t)).
\]
The flow identity
\(\partial_t\Phi^{T,t}(y)=-D_y\Phi^{T,t}(y)f(t,y)\) gives
\[
 \frac{d}{dt}\Phi^{T,t}(u(t))=M(T,t)h(t).
\]
Integrating from \(t_n\) to \(T\) yields the nonlinear
variation-of-constants formula
\begin{equation}\label{eq:nonlinear-defect-error-sep-23}
 x(T)-u(T)=-\int_{t_n}^T M(T,t)h(t)\,dt.
\end{equation}

Set \(\varphi(t)=M(T,t)h(t)\). By
\eqref{eq:defect-zeros-sep-23}, every quadrature sample vanishes, so
\[
 x(T)-u(T)=-E_{n,\tau}[\varphi].
\]
For sufficiently smooth \(f\), on the local branch of collocation
solutions for small \(\tau\), the required derivatives of
\(\varphi\) remain bounded. The quadrature error estimate therefore
gives
\begin{equation}\label{eq:collocation-local-order-sep-23}
 \boxed{x(t_n+\tau)-u(t_n+\tau)=O(\tau^{p+2}).}
\end{equation}
This proves the order \(p+1\) conclusion stated on September 21.
The extra accuracy at the endpoint comes from cancellation in the
integral of the defect; merely bounding \(\|h\|\) would miss that
cancellation.

\subsection{Gauss collocation}

Gaussian quadrature with \(s\) nodes has degree of precision
\(2s-1\), the largest possible value. Its nodes are the roots of the
shifted Legendre polynomial of degree \(s\):
\[
 L_s(\theta)=P_s(2\theta-1),\qquad 0< c_1<\cdots<c_s<1,
\]
where \(P_s\) is the Legendre polynomial on \([-1,1]\).
Equivalently, the monic nodal polynomial
\(\pi_s(\theta)=\prod_{j=1}^s(\theta-c_j)\) is orthogonal to
\(\mathcal P_{s-1}\) on \([0,1]\):
\[
 \int_0^1\pi_s(\theta)q(\theta)\,d\theta=0,
 \qquad q\in\mathcal P_{s-1}.
\]

To see the resulting precision, divide any
\(v\in\mathcal P_{2s-1}\) as
\(v=\pi_s q+r\), with \(q,r\in\mathcal P_{s-1}\).
The integral of \(\pi_s q\) vanishes by orthogonality, and all its
quadrature samples vanish as well. The interpolatory rule integrates
\(r\) exactly. Hence it integrates \(v\) exactly.

The nodes and weights satisfy
\begin{equation}\label{eq:gauss-node-properties-sep-23}
 c_j=1-c_{s+1-j},\qquad
 b_j=b_{s+1-j}>0.
\end{equation}
Positivity follows directly from exactness applied to
\(\ell_j^2\), whose degree is \(2s-2\):
\[
 b_j=\sum_{i=1}^s b_i\ell_j(c_i)^2
     =\int_0^1\ell_j(\theta)^2\,d\theta>0.
\]

Collocation at these nodes gives the \emph{Gauss RK method} with
\(s\) stages. The quadrature precision is \(2s-1\), so
\eqref{eq:collocation-local-order-sep-23} gives
\[
 \text{RK order}=2s,
 \qquad \text{local error}=O(\tau^{2s+1}).
\]

\subsubsection{One stage: implicit midpoint}

For \(s=1\), \(c_1=1/2\), \(b_1=1\), and \(a_{11}=1/2\):
\[
 \begin{array}{c|c}\tfrac12&\tfrac12\\ \hline &1\end{array}.
\]
Thus
\[
 k_1=f\left(t_n+\frac\tau2,x_n+\frac\tau2k_1\right),
 \qquad x_{n+1}=x_n+\tau k_1.
\]
Equivalently,
\[
 x_{n+1}=x_n+\tau f\left(t_n+\frac\tau2,
                              \frac{x_n+x_{n+1}}2\right).
\]
This method has order two. Its stability function is
\[
 R_{\mathrm{G},1}(z)=\frac{1+z/2}{1-z/2}.
\]
It agrees with the trapezoidal stability function on the linear test
equation, although the two methods differ for general nonlinear \(f\).

\subsubsection{Two-stage Gauss method}

The roots of \(L_2(\theta)=6\theta^2-6\theta+1\) are
\[
 c_1=\frac12-\frac{\sqrt3}{6},\qquad
 c_2=\frac12+\frac{\sqrt3}{6}.
\]
Integrating their Lagrange basis functions gives
\begin{equation}\label{eq:gauss-two-tableau-sep-23}
 \begin{array}{c|cc}
  \tfrac12-\tfrac{\sqrt3}{6}
    &\tfrac14&\tfrac14-\tfrac{\sqrt3}{6}\\[4pt]
  \tfrac12+\tfrac{\sqrt3}{6}
    &\tfrac14+\tfrac{\sqrt3}{6}&\tfrac14\\[4pt] \hline
    &\tfrac12&\tfrac12
 \end{array}.
\end{equation}
In particular, the upper-right entry has a minus sign, so that the
first row sums to \(c_1\). The stage equations are
\begin{align*}
 k_1&=f\left(t_n+c_1\tau,
       x_n+\tau\left[\frac14k_1+
                    \left(\frac14-\frac{\sqrt3}{6}\right)k_2\right]\right),\\*
 k_2&=f\left(t_n+c_2\tau,
       x_n+\tau\left[\left(\frac14+\frac{\sqrt3}{6}\right)k_1+
                    \frac14k_2\right]\right),\\*
 x_{n+1}&=x_n+\frac\tau2(k_1+k_2).
\end{align*}
Both slopes must be solved for together. The method has order four,
and substitution into \eqref{eq:rk-stability-function-sep-21} gives
\begin{equation}\label{eq:gauss-two-stability-sep-23}
 R_{\mathrm{G},2}(z)
  =\frac{1+z/2+z^2/12}{1-z/2+z^2/12}.
\end{equation}

\subsection{Quadratic first integrals}

For the autonomous equation \(x'=f(x)\), a function \(I(x)\) is a
\emph{first integral} if it is constant along every solution.
Consider a quadratic function
\[
 I(x)=x^{\mathsf T}Qx,\qquad Q^{\mathsf T}=Q.
\]
Positive definiteness is not needed for conservation. Since
\[
 \frac{d}{dt}I(x(t))=2x(t)^{\mathsf T}Qf(x(t)),
\]
the condition for a quadratic first integral is
\begin{equation}\label{eq:quadratic-invariant-condition-sep-23}
 x^{\mathsf T}Qf(x)=0\qquad\text{for every }x.
\end{equation}

\begin{theorem}[Quadratic invariants of Gauss methods]
 Every Gauss collocation method preserves each quadratic first
 integral exactly, provided its stage equations are solved exactly.
\end{theorem}

\begin{proof}
 Let \(U(\theta)=u(t_n+\tau\theta)\) and \(U_j=U(c_j)\).
 The polynomial \(U\) has degree at most \(s\), and
 \[
  U'(c_j)=\tau f(U_j).
 \]
 Therefore \(2U^{\mathsf T}QU'\) has degree at most \(2s-1\)
 and Gaussian quadrature is exact for it. Consequently,
 \begin{align*}
  I(x_{n+1})-I(x_n)
   &=\int_0^1 2U(\theta)^{\mathsf T}QU'(\theta)\,d\theta\\
   &=2\tau\sum_{j=1}^s b_j U_j^{\mathsf T}Qf(U_j)=0.
 \end{align*}
 The last equality uses \eqref{eq:quadratic-invariant-condition-sep-23}
 at each stage value.
\end{proof}

For example, if \(f(x)=Ax\) and \(A^{\mathsf T}Q+QA=0\), the
quadratic form is preserved by the exact solution and by every Gauss
method. Taking \(Q=I\) and \(A^{\mathsf T}=-A\) recovers norm
preservation for the oscillator from September 18.
The argument relies on the invariant being quadratic; it does not
give exact preservation of a general nonlinear first integral.

\subsection{An energy identity for the scalar test equation}

Apply Gauss collocation to \(x'=\lambda x\), with
\(z=\tau\lambda\). Now \(U\) is complex-valued and
\(U'(c_j)=zU_j\). Since the derivative of \(|U|^2\) is a real
polynomial of degree at most \(2s-1\), exact quadrature gives
\begin{align}
 |x_{n+1}|^2-|x_n|^2
  &=2\int_0^1\operatorname{Re}
       \bigl(\overline{U(\theta)}U'(\theta)\bigr)\,d\theta\notag\\
  &=2\operatorname{Re}z\sum_{j=1}^s b_j|U_j|^2.
       \label{eq:gauss-energy-identity-sep-23}
\end{align}
For \(\operatorname{Re}z<0\), positivity of the weights implies
\(|x_{n+1}|\leq|x_n|\).

This also proves unique solvability in the open left half-plane.
Indeed, for a homogeneous stage system take \(x_n=0\). The left
side is then nonnegative and the right side nonpositive, so every
\(U_j=0\). Since \(U(0)=0\) and the Gauss nodes are distinct and
interior, \(U\) has \(s+1\) distinct zeros and must vanish
identically. Thus the stage matrix has trivial kernel.

For \(x_n\neq0\), the sum in
\eqref{eq:gauss-energy-identity-sep-23} is strictly positive: if all
\(U_j=0\), then all \(U'(c_j)=0\), so \(U'\equiv0\) and
\(U\equiv0\), a contradiction. Therefore
\[
 \operatorname{Re}z<0\quad\Longrightarrow\quad |R(z)|<1.
\]
Every Gauss RK method is consequently \(A\)-stable. On the
imaginary axis, the same identity gives \(|R(z)|=1\) whenever the
stages are defined, explaining the absence of artificial damping for
oscillatory modes. We continue this comparison with Radau methods
in the next lecture.
